\documentclass[11pt,a4paper]{article}
\usepackage{fontspec}
\usepackage[margin=1in]{geometry}
\usepackage[unicode,hidelinks]{hyperref}
\setlength{\parindent}{0pt}
\setlength{\parskip}{.6em}
\emergencystretch=3em
\newcommand{\volumeentry}[3]{\par\begingroup\parskip=.7em
  \noindent\href{QRAIF_Complete_Collection.pdf\#page=#3}{\textbf{#1}\quad #2\nobreak\dotfill\nobreak\mbox{#3}}\par\endgroup}
\hypersetup{pdftitle={Quantum Reference Algebra for Information Flow},pdfauthor={Sidong Liu}}
% COVER SUBTITLE: World structure and physical control
\begin{document}
\begin{titlepage}
\centering
\vspace*{1.6cm}
{\Huge\bfseries Quantum Reference Algebra for Information Flow\par}
\vspace{1cm}
{\Large World structure and physical control\par}
\vfill
{\Large Complete Collected Volume\par}
\vspace{1.6cm}
{\Large Sidong Liu, PhD\par}
\vspace{.3cm}
{iBioStratix Ltd\par}
\href{mailto:sidongliu@hotmail.com}{\texttt{sidongliu@hotmail.com}}\par
\vspace{1cm}
{\large October 2026\par}
\end{titlepage}
\setcounter{page}{2}
\section*{Contents}
The numbers below refer to pages in this PDF. Each paper retains its own printed pagination, main argument and references. The technical appendices of Q-RAIF C are distributed as separate linked supplements; the other papers retain their appendices.
\volumeentry{}{Introduction and reading map}{3}
\volumeentry{QR-A}{Lorentzian Read-out from Local Unitary Propagation: Clifford Relations, Causal Cones and Spectral Mass}{4}
\volumeentry{QR-B}{Thermodynamic Stability Constraints on Persistent Open Quantum Subsystems: Control Algebra, Physical Actuation and Energy Cost}{28}
\volumeentry{QR-C}{Representing Control in the World: Retained History, Physical Sources and Renewal in Q-RAIF}{50}
\clearpage
\section*{Introduction}\label{collection:intro}
Why should the algebra through which a system acts be related to the structure through which its world propagates? Q-RAIF studies this interface directly. A positive control cost, an information metric and a Lorentzian form are distinct objects; a bridge has to follow the actual generators and their representation.

A derives the common quadratic shell and Clifford relations under its local propagation assumptions. B shows when a centred field-pulse cost is independent of the unknown carrier state, and when generator alignment relates that cost to the propagation form. C contrasts control discarded into a continuous-spectrum port with control retained in the closed propagation. The former restricts the exact multiplicative domain; the latter permits an exact represented algebra and a conditional Lorentzian extension.

Finite field actuators, material sources, clocks and receiver renewal provide physical witnesses and resource tests for these interfaces. They do not prove arbitrary world--control alignment or indefinitely cyclic operation. Mass--Energy, Null Cluster and No View remain separate companion papers: positive spectral readout, a matched null-access task and conservation-constrained source selection complete the six-paper responsibility without changing this volume's three-paper membership.

\section*{Reading map}
\begin{description}
\item[QR-A] Local propagation and the single quadratic shell constrain Clifford and Lorentzian readout. A common cone still needs the stated speed-alignment dynamics.
\item[QR-B] A centred physical coupling gives a state-independent cost criterion and a conditional bridge to the propagation form. Alignment is an additional physical question.
\item[QR-C] Discarded control ports and closed propagation give different exact representation results. Finite source, clock and renewal witnesses test their physical implementation.
\end{description}
\end{document}
